% !TeX root = ../main.tex
% Fonte: Lezione 9, pp. 1--4.

\section{Estrazione e completamento di basi}

Da una famiglia generatrice si arriva a una base eliminando vettori ridondanti; da una famiglia indipendente, aggiungendo vettori fino a generare lo spazio.

\begin{corollario}
Da ogni famiglia finita di generatori di uno spazio vettoriale si può estrarre una base.
\end{corollario}
\begin{proof}
\textbf{Eliminazione di un vettore dipendente.} Se $v_1,\ldots,v_m$ sono già indipendenti, formano una base. Altrimenti esiste una relazione $\sum_i\lambda_iv_i=0_V$ con $\lambda_j\ne0_K$ per un indice $j$, e quindi
\[
  v_j=-\sum_{i\ne j}\frac{\lambda_i}{\lambda_j}v_i.
\]

\textbf{Generazione invariata.} Per $v=\sum_i a_iv_i$, sostituendo l'espressione di $v_j$ si ottiene
\[
  v=\sum_{i\ne j}\left(a_i-a_j\frac{\lambda_i}{\lambda_j}\right)v_i.
\]
Togliere $v_j$ non modifica dunque il sottospazio generato.

\textbf{Terminazione.} Si ripete l'eliminazione finché la famiglia è dipendente. La cardinalità diminuisce a ogni passo e il sottospazio generato resta $V$: dopo un numero finito di passi rimane una famiglia libera di generatori, cioè una base. Se $V=\{0_V\}$ può rimanere la famiglia vuota.
\end{proof}

\begin{corollario}
Ogni famiglia linearmente indipendente in uno spazio vettoriale finitamente generato si può completare a una base.
\end{corollario}
\begin{proof}
Sia $V$ generato da $m$ vettori. Ogni famiglia indipendente in $V$ ha al più $m$ elementi: altrimenti se ne potrebbero scegliere $m+1$ contro il teorema dello scambio. Possiamo quindi scrivere la famiglia iniziale come $w_1,\ldots,w_s$.

\textbf{Scelta del nuovo vettore.} Se la famiglia genera $V$, è già una base. Altrimenti scegliamo $w_{s+1}\in V\setminus\langle w_1,\ldots,w_s\rangle$.

\textbf{Conservazione dell'indipendenza.} In una relazione
\[
  \lambda_1w_1+\cdots+\lambda_sw_s+\lambda_{s+1}w_{s+1}=0_V,
\]
un eventuale $\lambda_{s+1}\ne0_K$ permetterebbe di esprimere $w_{s+1}$ mediante i precedenti, contro la scelta. Dunque $\lambda_{s+1}=0_K$; l'indipendenza iniziale annulla anche gli altri coefficienti.

\textbf{Terminazione.} Finché manca la generazione, aggiungiamo un vettore esterno al sottospazio generato. L'indipendenza si conserva, ma lo scambio vieta di superare $m$ vettori. Dopo al più $m-s$ aggiunte il procedimento termina con una base di $V$.
\end{proof}

\section{Basi di sottospazi di \texorpdfstring{$\R^4$}{R4}}

\subsection{Da equazioni cartesiane a una base}

\begin{esercizio}
Trovare una base del sottospazio
\[
  U=\left\{(x_1,x_2,x_3,x_4)\in\R^4:
  \begin{cases}
    2x_1-x_2-3x_4=0,\\
    2x_1+x_3+x_4=0
  \end{cases}\right\}.
\]
\end{esercizio}

\begin{svolgimento}
\textbf{Parametrizzazione.} Le equazioni danno $x_2=2x_1-3x_4$ e $x_3=-2x_1-x_4$. Ponendo $x_1=a_1$ e $x_4=a_2$ otteniamo
\begin{align*}
  (x_1,x_2,x_3,x_4)
  &=(a_1,2a_1-3a_2,-2a_1-a_2,a_2)\\
  &=a_1(1,2,-2,0)+a_2(0,-3,-1,1).
\end{align*}

\textbf{Generazione.} Ogni elemento di $U$ ha questa forma e, viceversa, ogni combinazione dei vettori
\[
  u_1=(1,2,-2,0),\qquad u_2=(0,-3,-1,1)
\]
soddisfa le due equazioni. Pertanto $U=\langle u_1,u_2\rangle$.

\textbf{Indipendenza.} Se $a_1u_1+a_2u_2=0$, la prima coordinata impone $a_1=0$ e la quarta $a_2=0$. Quindi $(u_1,u_2)$ è una base e $\dim U=2$.
\end{svolgimento}

\subsection{Da una famiglia di generatori a una base}

\begin{esercizio}
Si consideri il sottospazio di $\R^4$ generato da
\[
  v_1=(1,-1,0,2),\qquad v_2=(-1,-3,2,0),\qquad v_3=(2,0,-1,3).
\]
Verificare che ha dimensione $2$ e trovarne una base.
\end{esercizio}

\begin{svolgimento}
\textbf{Relazioni tra i generatori.} L'uguaglianza $\lambda_1v_1+\lambda_2v_2+\lambda_3v_3=0$ equivale a
\[
  \begin{cases}
    \lambda_1-\lambda_2+2\lambda_3=0,\\
    -\lambda_1-3\lambda_2=0,\\
    2\lambda_2-\lambda_3=0,\\
    2\lambda_1+3\lambda_3=0.
  \end{cases}
\]
La seconda e la terza equazione danno $\lambda_1=-3\lambda_2$ e $\lambda_3=2\lambda_2$; le altre due risultano identicamente soddisfatte. Per $\lambda_2=1$ si trova
\[
  -3v_1+v_2+2v_3=0,\qquad v_2=3v_1-2v_3.
\]

\textbf{Estrazione della base.} Si elimina $v_2$ senza modificare il sottospazio generato. I vettori $v_1,v_3$ sono indipendenti: in $av_1+bv_3=0$, la seconda coordinata dà $a=0$ e la terza $b=0$. Dunque $(v_1,v_3)$ è una base e la dimensione è $2$.
\end{svolgimento}

\subsection{Parametri liberi e vettori di base}

\begin{esercizio}
Trovare una base del sottospazio
\[
  W=\left\{(x_1,x_2,x_3,x_4)\in\R^4:
  \begin{cases}
    x_1+2x_3+x_4=0,\\
    x_2-x_4=0
  \end{cases}\right\}.
\]
\end{esercizio}

\begin{svolgimento}
\textbf{Scelta dei parametri.} Si ha $x_1=-2x_3-x_4$ e $x_2=x_4$. Le coordinate $x_3,x_4$ sono libere: ponendo $x_3=a_1$ e $x_4=a_2$,
\begin{align*}
  (x_1,x_2,x_3,x_4)
  &=(-2a_1-a_2,a_2,a_1,a_2)\\
  &=a_1(-2,0,1,0)+a_2(-1,1,0,1).
\end{align*}

\textbf{Generazione e indipendenza.} I vettori corrispondono ai parametri $(1,0)$ e $(0,1)$. Tutte le loro combinazioni soddisfano le equazioni e ogni elemento di $W$ ha questa forma. In una combinazione nulla, la terza coordinata impone $a_1=0$ e la quarta $a_2=0$. Pertanto
\[
  \bigl((-2,0,1,0),\,(-1,1,0,1)\bigr)
\]
è una base di $W$ e $\dim W=2$.
\end{svolgimento}

